\origpage{278}{ROBERTSON AND SEYMOUR}
\hspace*{12pt}(3) \textit{Every component of $G$ realizes a unique vertex of $V(G^h)$.}

By (1) and (2), it suffices to show that if $x \ne y$ and $P$ is a path of $G$ between $x \in V(\sigma_m(u))$ and $y \in V(\sigma_n(v))$ then $u = v$; and it clearly suffices to prove this when no internal vertex of $P$ is in $V(Q)$. We assume $m \le n$ without loss of generality, and we proceed by induction on $n - m$. Suppose that $e \in E(Q)$ for some edge $e$ of $P$. Then both ends of $e$ are in $V(Q)$, and so $E(P) = \{e\}$. Choose $l$ with $1 \le l \le r_h$ and $w \in V(G^h_l)$ such that $e \in E(\sigma_l(w))$. By (2), $w = u$ and $w = v$, and so $u = v$, as required. We may assume then that every edge of $P$ is in $E(R)$. Since $E(P) \ne \emptyset$ and $E(P) \subseteq E(R)$, $P$ is a subpath of $P_p$ for some $p$ with $1 \le p \le k$. By (14.1), no vertex of $P$ except $x$ is in $V(G^i_{j(m)})$ (and so $m \ne n$), and similarly, no vertex of $P$ except $y$ is in $V(G^i_{j(n)})$. In particular, $x = \tau^i_{j(m)}(p+k)$ and $y = \tau^i_{j(n)}(p)$. If $n \ge m + 2$ then by (14.1) some vertex of $P$ is in $V(G^i_{j(m+1)})$, and the last such vertex is in $Y \subseteq V(Q)$. Since no internal vertex of $P$ is in $V(Q)$, we deduce that one of $x$, $y$ is in $V(G^i_{j(m+1)})$, and the result follows from (2) and our inductive hypothesis. We may assume then that $n = m + 1$. But $x = \tau^i_{j(m)}(p+k)$ and so $u = \tau^h_m(p+k)$. Similarly, $v = \tau^h_n(p)$. But $\tau^h_m(p+k) = \tau^h_n(p)$ since $n = m + 1$, and so $u = v$, as required.

\hspace*{12pt}(4) \textit{Every vertex of $V(G^h)$ is realized by a unique component of $G$.}

Certainly every vertex of $V(G^h)$ is realized by some component of $G$. Let $v \in V(G^h)$, and let $x \in V(\sigma_m(v))$, $y \in V(\sigma_n(v))$. We wish to prove that $x$ and $y$ are in the same component of $G$. If $m = n$ this is true since $\sigma_m(v)$ is connected, and so we may assume $n > m$. Suppose that $n = m + 1$. Then $v \in V(G^h_m) \cap V(G^h_n)$, and so $v \in Y^h_m$. Choose $p$ with $1 \le p \le k$ such that $v = \tau^h_m(p+k) = \tau^h_n(p)$. Then $x$ is the same component of $G$ as $\tau^i_{j(m)}(p+k)$, and $y$ is in the same component of $G$ as $\tau^i_{j(n)}(p)$. But $\tau^i_{j(m)}(p+k)$ and $\tau^i_{j(n)}(p)$ are in the same component of $R$ and hence of $G$, by (14.1), since $n = m + 1$. The result is therefore true if $n = m + 1$. Now suppose $n > m + 1$. Then $v \in V(G^h_l)$ for $m < l < n$, and we may choose $z_l \in V(G^i_{j(l)})$ ($m \le l \le n$) such that $z_l \in V(\sigma_l(v))$, and $z_m = x$, $z_n = y$. But for $m \le l < n$, $z_l$ and $z_{l+1}$ are in the same component of $G$, by our observation above. Hence $x$ and $y$ are in the same component of $G$, as required.

Now to complete the proof of (14.2), we define $\sigma$ by (i) if $v \in V(G^h)$, $\sigma(v)$ is the component of $G$ which realizes $v$ (ii) if $e \in E(G^h)$, $\sigma(e)$ is $\sigma_n(e)$, where $1 \le n \le r_h$ and $e$ is an edge of $G^h_n$.

It is easy to verify that $\sigma$ is a simulation of $(G^h, T^h)$ in $(G^i, T^i)$, as required. This completes the proof of (14.2). $\square$

Now we need a mild generalization of $k$-chains. We say that $(G, T)$ is a
